\section{西晋接收与三国叙事的后效}

二八〇年以后，西晋最先接收的是可在洛阳完成的礼仪：孙皓出降、受封和安置，
使军事终点进入新朝的名义秩序。更慢的是江东的州郡、港口、旧军和县廷。哪些
吴吏留任，军队如何遣散，税粮向何处转运，地方何时改用晋朝的法律与任命语言，
都不可能由一份降表同时完成。《晋书》《三国志·孙皓传》及《资治通鉴》能
确认伐吴与出降的骨架，不能为每座城、每个家庭补出接收日程。

这也改变我们读魏晋禅代的方式。禅诏能排列名义交接，无法替地方解释它是否
已转化为保护、征发或新的官司程序。西晋继承魏的近畿官署、蜀汉的山道经验和
孙吴的水路网络，继承的同时也是三种不等速的成本。泰始律提供一个可援引的
共同规范，并不证明州县已同速执行；法律文本与地方实践之间仍隔着书记、官吏、
家族和道路。

三国的后效还在史书中延续。陈寿以魏纪年组织《三国志》，使魏的中枢时间成为
全书骨架；蜀书、吴书和裴松之所引杂史则保留了可抵抗单线叙事的材料。裴注的
增补并不自动更接近现场，它包含不同成书年代、家传立场和地域记忆，须逐条
辨认。把《华阳国志》、吴简、石经残片和考古城址带入阅读，不是为了拼出一部
无缝全史，而是为了让皇帝、道路、县廷、山地与海峡各有不同的可见尺度。

因此“统一”既是一个必要的编年终点，也不应成为叙事的封口。它结束了三个
最高政权并立，却没有消除地方社会的差异，更没有让后世史家的取舍失去作用。
读者回看曹魏、蜀汉和孙吴时，应保留《三国志》简短动作如何推进选择与后果的
力量，同时拒绝以它的魏中心框架替所有未被充分保存的人和地方说话。

\subsection{接收蜀地，预演接收江东}

灭蜀后钟会集团的动荡说明，得到成都并不等于已得到蜀地军人和命令网络。魏军
内部、降众和姜维的后续行动使接收者必须立刻处理互不相同的忠诚与恐惧。司马
氏由此得到一项并不轻松的经验：降表可以结束一座都城的最高政权，不能替新朝
安置军队、解释任命或修复道路。伐吴前对益州造船和荆州部署的经营，正是在
这种经验之上把不同区域重新放入同一场准备。

西晋史家如何书写这种过程，同样构成后效。《晋书》的成书晚于事件甚久，保存
了新朝回望统一的语言；它可与三国志、通鉴互校次序，不能独自为军议和地方
社会提供透明窗口。后世“天下一统”的叙述很容易把二八〇年之前的多种地方
时间压平。把吴简、山城考古、石经残片和不同政权的纪传置回阅读，既不否认
统一的政治意义，也拒绝让统一替未被保存的经验消音。
